\documentclass[a4paper, 12pt]{article}

% Essential Packages
\usepackage{amsmath}
\usepackage{amssymb}
\usepackage[top=25mm, bottom=25mm, left=25mm, right=25mm]{geometry}
\usepackage{pgfplots}
\usepackage[inline]{enumitem}
\usepackage{mathtools}
\usepackage{tikz}
\usepackage{fancyhdr}
\usepackage{setspace}
\usepackage{hyperref}

% Misc
\pgfplotsset{compat=1.18}
\usepgfplotslibrary{polar}
\usepgfplotslibrary{fillbetween}
\setlength{\parskip}{3pt}
\setstretch{1.15}
\renewcommand{\headrulewidth}{0pt}
\renewcommand{\footrulewidth}{0pt}
\fancyhf{}

\begin{document}
\pagestyle{empty}

\begin{center}
\textbf{2024-2025 Fall \\MAT123 Midterm\\14/11/2025}
\end{center}

\begin{enumerate}[leftmargin=*,label=\textbf{\arabic*.}]
\item Given that
\[
f(x) =
\begin{cases}
e^{-1/x^2}, & \text{if } x\neq0 \\[1em]
0, & \text{otherwise}
\end{cases},
\]
find all critical and singular points of $f(x)$. Determining the intervals where $f(x)$ increases and where it decreases, find local extreme values.

\item Determine whether each of the statement below is true or false, providing short explanations if it is true and counterexamples if it is false.

\begin{enumerate}[leftmargin=*,label=\textbf{\alph*.}]
\item Every increasing function has inverse.
\item If a function is continuous on an interval $I$, then it must be differentiable on $I$.
\item If $f'(c)=0$ for some point $c$ in the domain of a function $f$, then $f(c)$ is either a local minimum or a local maximum value of $f$.
\item If both $f'(x)$ and $f''(x)$ exist, with $f''(x)\neq0$ and $f'(x)/f''(x)>0$ on the open interval $(a,b)$, then $f(x)$ must be either increasing or decreasing on $(a,b)$.
\end{enumerate}

\item Compute the following limits.
\begin{enumerate}[leftmargin=*,label=\textbf{\alph*.}]
\item $\displaystyle\lim_{t\to\pi/2}\frac{\left|\cos t-\sin t\right|}{\left|\cos t\right|-\left|\sin t\right|}$
\item $\displaystyle\lim_{x\to\infty} x^{\sin\left(\frac1x\right)}$
\end{enumerate}

\item Find the point(s) on the curve $2\tan^{-1}(y/x)=\ln\left(x^2+y^2\right)$ at which the tangent line is horizontal.

\item A point P in the first quadrant of the Cartesian plane is moving along a circle centered at the origin. As the point $P$ moves, the $x$-intercept of the tangent to the circle at $P$ moves on the $x$-axis at the rate of $1$ unit per second. At what rate is the point $P$ moving when the tangent line has slope $-1$?
\end{enumerate}

\newpage
\pagestyle{fancy}
\renewcommand{\headrulewidth}{1pt}
\fancyhead[R]{\textit{2025-2026 Fall Midterm}}
\fancyhead[L]{\textit{Hacettepe MAT123 Exam Solutions Version 2}}

%Last update: 03/09/2026 ??:??

\begin{enumerate}[leftmargin=*,label=\textbf{\arabic*.}]
\item For $x \ne 0$, $f'(x) = \dfrac{2}{x^3} e^{-1/x^2}$ by the chain rule. Also,
\[f'(0) = \lim_{x \to 0} \frac{e^{-1/x^2}}{x} \overset{\text{L'H.}}{=} \lim_{x\to0}\frac{1/x}{e^{1/x}}=\lim_{t\to\pm\infty}\frac t{e^t}\overset{\text{L'H.}}{=}\lim_{t\to\pm\infty}\frac1{e^t}=0.\]

Hence,
\[f'(x) = \begin{cases} \dfrac{2}{x^3} e^{-1/x^2}, & \text{if } x \ne 0 \\ 0, & \text{if } x = 0 \end{cases}.\]
\[\boxed{
\begin{gathered}
\text{The only critical point is } x=0. \text{ There are no singular points.}\\
f\text{ is increasing on }(-\infty,0)\text{ and decreasing on }(0,\infty).\\
f(0)=0\text{ is the only local (and also absolute) minimum value of }f.
\end{gathered}
}\]

\item 
\begin{enumerate}[leftmargin=*,label=\textbf{\alph*.}]
\item $\boxed{\text{TRUE.}}$ For $x_1, x_2$ in the domain of an increasing function $f$ with $x_1 \ne x_2$, if \linebreak $x_1 < x_2$ then $f(x_1) < f(x_2)$, and if $x_1 > x_2$ then $f(x_1) > f(x_2)$. Thus, $f(x_1) \ne f(x_2)$, so $f$ is one-to-one and has an inverse.

\item $\boxed{\text{FALSE.}}$ $f(x) = |x|$ is continuous on $(-\infty, \infty)$ but not differentiable at $x = 0$.

\item $\boxed{\text{FALSE.}}$ Consider $f(x) = x^3$, whose critical point is $x = 0$ ($f'(0) = 0$). However, $f(0) = 0$ is neither a local maximum nor a local minimum value.

\item $\boxed{\text{TRUE.}}$ Since $f'(x)/f''(x) > 0$, $f'(x)$ and $f''(x)$ must have the same sign and $f'(x) \ne 0$ on $(a, b)$. By the Intermediate Value Theorem, $f'(x)$ cannot change sign on $(a, b)$ without crossing zero. Thus, either $f'(x) > 0$ for all $x \in (a, b)$ or $f'(x) < 0$ for all $x \in (a, b)$, meaning $f(x)$ is strictly increasing or strictly decreasing on $(a, b)$.
\end{enumerate}

\item 
\begin{enumerate}[leftmargin=*,label=\textbf{\alph*.}]
\item Note that for $t$ close to $\frac{\pi}{2}$, $\cos t - \sin t < 0$, $\left|\cos t\right| = \cos t$, and $\left|\sin t\right| = \sin t$.
\[\lim_{t \to \frac{\pi}{2}^-} \frac{\left|\cos t - \sin t\right|}{\left|\cos t\right| - \left|\sin t\right|} = \lim_{t \to \frac{\pi}{2}^-} \frac{\sin t - \cos t}{\cos t - \sin t} = -1\]
\[\lim_{t \to \frac{\pi}{2}^+} \frac{\left|\cos t - \sin t\right|}{\left|\cos t\right| - \left|\sin t\right|} = \lim_{t \to \frac{\pi}{2}^+} \frac{\sin t - \cos t}{-\cos t - \sin t} = \frac{1 - 0}{-0 - 1} = -1\]

Since both one-sided limits equal $-1$, the limit is $\boxed{-1}$.

\item Let $y = x^{\sin(1/x)}$. Then $\ln y = \sin(1/x) \ln x$.
\[\lim_{x \to \infty} \sin\left(\frac{1}{x}\right) \ln x = \lim_{x \to \infty} \frac{\ln x}{\csc(1/x)} \quad \left[\frac{\infty}{\infty}\right]\]

Applying L'Hôpital's Rule:
\[= \lim_{x \to \infty} \frac{1/x}{-\csc(1/x) \cot(1/x) (-1/x^2)} = \lim_{x \to \infty} \tan\left(\frac{1}{x}\right) \frac{\sin(1/x)}{1/x} = 0 \cdot 1 = 0\]
Thus, $\displaystyle\lim_{x \to \infty} x^{\sin(1/x)} = e^0 = \boxed1$.
\end{enumerate}

\item Differentiating both sides implicitly with respect to $x$,
\[2 \frac{\frac{d}{dx}(y/x)}{1 + (y/x)^2} = \frac{\frac{d}{dx}(x^2 + y^2)}{x^2 + y^2},\]
\[2 \frac{\dfrac{y'x - y}{x^2}}{\dfrac{x^2 + y^2}{x^2}} = \frac{2x + 2yy'}{x^2 + y^2}.\]

Simplifying gives $y'x - y = x + yy'$, so $y'(x - y) = x + y$, yielding
\[y' = \frac{x + y}{x - y}.\]

For horizontal tangents, we set $y' = 0$, which implies $y = -x$. Substituting $y = -x$ into the original equation,
\[2 \tan^{-1}(-1) = \ln(2x^2) \implies 2 \left(-\frac{\pi}{4}\right) = \ln(2x^2) \implies 2x^2 = e^{-\pi/2},\]
\[x = \pm \frac{1}{\sqrt{2}} e^{-\pi/4}.\]

Since $y = -x$, the points are
\[\boxed{\left(\frac{1}{\sqrt{2}} e^{-\pi/4}, -\frac{1}{\sqrt{2}} e^{-\pi/4}\right) \quad \text{and} \quad \left(-\frac{1}{\sqrt{2}} e^{-\pi/4}, \frac{1}{\sqrt{2}} e^{-\pi/4}\right)}.\]

\item Let the circle have radius $r$, so its equation is $x^2 + y^2 = r^2$. A point $P$ in the first quadrant can be parametrized as $P(r \cos \theta, r \sin \theta)$ for $\theta \in (0, \pi/2)$. The slope of the tangent line at $P$ is $m =-\frac xy= -\cot \theta$. The equation of the tangent line at $P$ is
\[y - r \sin \theta = -\cot \theta (x - r \cos \theta) \qquad[y-y_0=m(x-x_0)].\]

Setting $y = 0$, the $x$-intercept is $x_0 = r \sec \theta$. We are given $\left.\dfrac{dx}{dt}\right|_{x=x_0} = 1$. Differentiating $x = r \sec \theta$ with respect to $t$,
\[\frac{dx}{dt} = r \sec \theta \tan \theta \frac{d\theta}{dt} = 1 \implies \frac{d\theta}{dt} = \frac{1}{r \sec \theta \tan \theta}.\]

The speed of point $P$ moving along the circle is $v = r \left|\frac{d\theta}{dt}\right| = \frac{1}{\sec \theta \tan \theta} = \frac{\cos^2 \theta}{\sin \theta}$.

When the tangent line has slope $-1$, $-\cot \theta = -1 \implies \theta = \frac{\pi}{4}$. At $\theta = \frac{\pi}{4}$:
\[v = \frac{\cos^2(\pi/4)}{\sin(\pi/4)} = \frac{(1/\sqrt{2})^2}{1/\sqrt{2}} = \boxed{\frac{1}{\sqrt{2}} \text{ units/sec}}.\]
\end{enumerate}
\end{document}